\subsection{APPLICATION: LOGARITHMIC CRITERIA FOR CONVERGENCE OF INTEGRALS}

By choosing the function g suitably one can deduce criteria for deciding whether the integral \(\int_{a}^{+\infty} f(t) dt\) of a function \(f \geqslant 0\) converges or is infinite from prop. 1 of V, p.~228, and cor. 1 thereto: it suffices to choose for g a function whose primitive is known. In particular, since \(x^{\mu}\) has primitive \(\frac{x^{\mu+1}}{\mu+1}\) when \(\mu \neq -1\) , and \(\log x\) when \(\mu = -1\) , we have the following criterion:

PROPOSITION 2 (``logarithmic criterion of order 0''). Let f be a regulated function \(\geqslant 0\) on an interval \([a, +\infty[\); if \(f(x) \preccurlyeq x^{\mu}\) for some \(\mu < -1\) then the integral \(\int_{a}^{+\infty} f(t) dt\) converges; if \(f(x) \succcurlyeq x^{\mu}\) for some \(\mu \geqslant -1\) then the integral \(\int_{a}^{+\infty} f(t) dt\) is infinite.

This criterion is not decisive when \(1 / x^{1 + \alpha} \ll f(x) \ll 1 / x\) for every exponent \(\alpha > 0\) , as, for example, when \(f(x) = 1 / x(\log x)^{\mu}\) ( \(\mu > 0\) ). But in this last case \(f\) has primitive \(\frac{1}{1 - \mu} (\log x)^{1 - \mu}\) if \(\mu \neq 1\) and \(\log \log x\) when \(\mu = 1\) . Thus:

PROPOSITION 3 (``logarithmic criterion of order 1''). Let f be a regulated function \(\geqslant 0\) on an interval \([a, +\infty[\); if \(f(x) \preccurlyeq 1/x(\log x)^{\mu}\) for some \(\mu > 1\) then the integral \(\int_{a}^{+\infty} f(t) dt\) converges; if \(f(x) \succcurlyeq 1/x(\log x)^{\mu}\) for some \(\mu \leqslant 1\) then the integral \(\int_{a}^{+\infty} f(t) dt\) is infinite.

Generally, for every integer \(n \geqslant 0\) , let us denote by \(l_{n}(x)\) the function defined inductively (for x large enough) by the relations \(l_{0}(x) = x\) , \(l_{n}(x) = \log(l_{n-1}(x))\) for \(n \geqslant 1\) ; one says that \(l_{n}(x)\) is the \(n^{th}\) iterated logarithm of x (cf.~Appendix). One verifies immediately that \(\frac{1}{1 - \mu}\left(l_{n}(x)\right)^{1 - \mu}\) is a primitive of

\[
\frac {1}{x l _ {1} (x) l _ {2} (x) \dots l _ {n - 1} (x) \left(l _ {n} (x)\right) ^ {\mu}}
\]

for \(\mu \neq 1\) , and \(l_{n+1}(x)\) is a primitive of \(\frac{1}{xl_1(x)l_2(x)\ldots l_{n-1}(x)l_n(x)}\) . Hence:

PROPOSITION 4 (``logarithmic criterion of order n''). Let f be a regulated function \(\geqslant0\) on an interval \([a,+\infty[\); if, for some \(\mu>1\), \(f(x)\preccurlyeq\frac{1}{xl_{1}(x)l_{2}(x)\ldots l_{n-1}(x)(l_{n}(x))^{\mu}}\) , then the integral \(\int_{a}^{+\infty}f(t)dt\) converges; if \(f(x)\succcurlyeq\frac{1}{xl_{1}(x)l_{2}(x)\ldots l_{n-1}(x)(l_{n}(x))^{\mu}}\) for some \(\mu\leqslant1\) , then the integral \(\int_{a}^{+\infty}f(t)dt\) is infinite.

Each logarithmic criterion is thus applicable to functions for which the criteria of lower order are not decisive (cf.~V, p.~264, exerc. 5 b) and V, p.~265, exerc. 8).

Because of its usefulness we translate the criterion of order 0 for integrals \(\int_{\alpha}^{a}f(t)dt\) where \(f\) is regulated and \(\geqslant 0\) on a non-compact interval \([\alpha ,a]\) :

PROPOSITION 5 (``logarithmic criterion of order 0''). If on a neighbourhood of \(\alpha\) one has \(f(x) \preccurlyeq 1/(x - \alpha)^{\mu}\) for some \(\mu < 1\) then the integral \(\int_{\alpha}^{a} f(t) dt\) converges; if \(f(x) \succcurlyeq 1/(x - \alpha)^{\mu}\) for some \(\mu \geqslant 1\) then the integral \(\int_{\alpha}^{a} f(t) dt\) is infinite.

We leave it to the reader to translate the logarithmic criterion of order \(n\) similarly.

The application of the logarithmic criteria is immediate if one knows how to obtain the principal part of f with respect to a comparison scale containing the functions that feature in these criteria: if \(f_{1}\) is the principal part, the integral \(\int_{\alpha}^{+\infty} f(t) dt\) converges or is infinite together with \(\int_{\alpha}^{+\infty} f_{1}(t) dt\) , and the logarithmic criteria apply immediately to this last integral.

Examples. 1) The function \(t^p (1 - t)^q\) is not bounded on {]}0, 1{[} when \(p < 0\) or \(q < 0\) ; by the logarithmic criterion of order 0 applied on a neighbourhood of the points 0 and 1 the integral \(\int_0^1 t^p (1 - t)^q dt\) converges if and only if \(p > -1\) and \(q > -1\) . If so, this integral is called the Eulerian integral of the first kind and is denoted by \(\mathbf{B}(p + 1,q + 1)\) (cf.~VII, p.~312).

\begin{enumerate}
\def\labelenumi{\arabic{enumi})}
\setcounter{enumi}{1}
\tightlist
\item
  Consider the integral \(\int_0^\infty t^{x - 1}e^{-t}dt\) . Since \(e^{-t}\sim 1\) on a neighbourhood of 0, one must have \(x > 0\) if this integral is to converge; this condition is also sufficient since on a neighbourhood of \(+\infty\) one has \(e^{-t}\ll t^{-\mu}\) for any \(\mu >0\) . When \(x > 0\) the integral is called the Eulerian integral of the second kind and is denoted by \(\Gamma (x)\) (cf.~VII, p.~311).
\end{enumerate}

